\section{Conceptual Framework：如何把本期内容压成一张心智地图}

苏煜在访谈结尾说，自己喜欢 build conceptual framework。这个表达也适合用来理解本期内容：Agent 不是一个孤立产品，而是一套概念连接。它从早期 AI 的“智能体”理想出发，经过符号逻辑、神经网络、语义解析，在 LLM 出现后重新获得统一接口；随后借助 coding、工具使用、world model 和 continual learning，开始进入生产系统。

\begin{importantbox}{一张口头心智地图}
本期可以压缩成四层：第一层是定义，Agent 是有边界、在环境中目标导向行动的实体；第二层是历史，Logical Agent、Neural Agent、Semantic Parsing、Language Agent 是同一问题的不同阶段；第三层是接口，language/coding/tool/API/GUI/CLI 正在收敛；第四层是生产化，可靠性、速度、成本、安全和持续学习决定它能否落地。
\end{importantbox}

\subsection{为什么“边界消融”不是口号}

边界消融的意思不是所有界面都会消失，而是同一个任务可以被不同接口表达。浏览器操作可以被 DOM 和脚本表达，桌面操作可以被 UI automation 表达，API 调用可以被代码封装，coding 又可以生成和改写这些接口。Agent 真正需要学习的是任务语义、环境状态和可行动作空间，而不是被某个入口形式锁死。

\subsection{从研究问题到产品问题}

研究阶段常问“这个 agent 能不能做”；产品阶段必须问“它能不能稳定、便宜、安全地做”。这也是为什么本期多次回到 reliability、speed、cost 和 forward deployment。OpenClaw Moment 让大家看见可能性，但生产系统需要让可能性穿过组织、权限、数据、遗留系统和用户习惯。

\begin{warningbox}{概念框架的误用}
概念框架不是万金油。把所有东西都放进一个框架，容易让差异消失。本期框架的价值在于解释趋势，但具体产品仍要回到任务、数据、用户、环境和成本。不同场景下，GUI、CLI、API、coding 的相对重要性会不同。
\end{warningbox}

\subsection{本章小结}

本期最强的内容不是某个单点预测，而是一个框架：Agent 是历史问题的新统一，language 是接口，coding 是数字世界的形式层，continual learning 是下一阶段能力积累路径，生产化则由可靠性和成本决定。

